\section{Resampling with the permutations left in the transform}
\label{sec:resampling}

The transform lengths needed for the integer product are pairwise coprime,
whereas the synthetic transforms use powers of two. Gaussian resampling
connects these two shapes. Its usual implementation sorts twice on each
coordinate. We retain both permutations in the transform formula and
compute that formula directly. This section proves the resulting interface,
including its tape cost when the number of coordinates grows.

For an integer $q\geq 1$, write
\[
 (F_q u)_k=\frac1q\sum_{j=0}^{q-1}u_j e^{-2\pi i jk/q},
 \qquad
 (F_q^+ u)_k=\frac1q\sum_{j=0}^{q-1}u_j e^{2\pi i jk/q}.
\]
All subscripts on a vector of length $q$ are taken modulo $q$.
The vector norm is the maximum modulus of its coordinates. All operator
norms are induced by this norm: $\|L\|=\sup_{\|u\|\leq1}\|Lu\|$.
Each normalized Fourier transform above is a contraction, since each
output averages inputs multiplied by complex numbers of modulus one. Coordinate
permutations preserve the norm.
For an integer precision $p\geq1$, put
\[
 \mathbb D_p=\{2^{-p}(a+ib):a,b\in\mathbb Z,\ a^2+b^2\leq 2^{2p}\}.
\]
For storage between completed scalar routines, put $w_{\rm c}=p+2$.
Represent $2^{-p}(a+ib)\in\mathbb D_p$ by a signed $w_{\rm c}$-bit
two's-complement word for $a$, followed by one for $b$, with each word
written most significant bit first. A coordinate list is the
concatenation of these fixed-length records in its stated order. Since
$|a|,|b|\leq2^p<2^{w_{\rm c}-1}$, this format includes both signs of
every permitted numerator, including $-2^p$ and $2^p$.
An approximation $\widetilde L:\mathbb D_p^a\to\mathbb D_p^b$
to a linear contraction $L:\mathbb C^a\to\mathbb C^b$ has error $E$
if $2^p\|\widetilde L u-Lu\|<E$ for every input in $\mathbb D_p^a$.
The approximation itself need not be linear. Composing such algorithms
adds their errors: at each step, insert the exact map at the computed
input and use its contraction bound on the accumulated input error.

We use the established $O(N\log N)$ integer multiplier of
Harvey and van der Hoeven~\cite[Theorem~1.1]{HarveyHoeven2021}
for the shorter integer records that occur below. We also use their
Lemma~2.12: for integer precision $p>100$ and integers $k>1$, $0\leq j<k$,
one can compute $\widetilde\zeta\in\mathbb D_p$ with
$2^p|\widetilde\zeta-e^{2\pi i j/k}|<2$ in time
$O(\max(p,\log k)^{1+\delta})$, for every fixed $0<\delta<1/8$.
Throughout this section, numbered citations to that paper refer to its
author-hosted 45-page manuscript.

\subsection{The one-dimensional interface}

\begin{lemma}[Resampling without executing the two permutations]
\label{lem:no-sort-resampling}
Fix $0<\delta<1/8$. Let $p>100$, $2\leq s<t<2^p$ and $2\leq\alpha<\sqrt p$
be integers, with $\gcd(s,t)=1$ and
\[
 \theta:=t/s-1>p/\alpha^4.
\]
Define coordinate permutations by
\[
 (P_su)_j=u_{tj},\qquad (P_tv)_k=v_{-sk}.
\]
There are linear contractions $A:\mathbb C^s\to\mathbb C^t$ and
$B_0:\mathbb C^t\to\mathbb C^s$ satisfying
\begin{equation}
 P_sF_s=2^{2\alpha^2}B_0P_tF_tA.
 \label{eq:no-sort-one-dimensional}
\end{equation}
There are fixed multitape algorithms for approximations
$\widetilde A:\mathbb D_p^s\to\mathbb D_p^t$ and
$\widetilde B_0:\mathbb D_p^t\to\mathbb D_p^s$, each with error less
than $p^2$ and time
\[
 O(tp^{3/2+\delta}\alpha).
\]
The constant is uniform in $p,s,t,\alpha$. Inputs and outputs use the
scalar component format above and occur in increasing coordinate order.
\end{lemma}

\begin{proof}
The Gaussian resampling argument uses two maps
$\mathsf S,\mathsf T:\mathbb C^s\to\mathbb C^t$. For $0\leq k<t$, they are
\[
 (\mathsf S u)_k=\frac1\alpha\sum_{j\in\mathbb Z}
 e^{-\pi\alpha^{-2}(j-sk/t)^2}u_j,
 \qquad
 (\mathsf T u)_k=\sum_{j\in\mathbb Z}
 e^{-\pi\alpha^2(tj/s-k)^2}u_j.
\]
Both sums converge absolutely. For $0\leq j<s$, define
\[
 q_j=\left\lfloor\frac{tj}{s}+\frac12\right\rfloor,
 \qquad \beta_j=\frac{tj}{s}-q_j.
\]
Thus $q_j$ is an integer nearest to $tj/s$, with ties rounded upward.
Let $C:\mathbb C^t\to\mathbb C^s$ select these coordinates, and let
$\mathsf D:\mathbb C^s\to\mathbb C^s$ be the diagonal map
\[
 (Cv)_j=v_{q_j},\qquad
 (\mathsf D u)_j=e^{\pi\alpha^2\beta_j^2}u_j.
\]
Set $N=C\mathsf T\mathsf D$. In $(Nu)_j$, the summand with index $j$
is exactly $u_j$: the factor in $\mathsf D$ cancels the Gaussian factor
at the selected coordinate $q_j$. The other summands form $((N-I)u)_j$.
The exact Gaussian identity and inverse estimate proved in
\cite[Theorem~4.2 and Lemma~4.6]{HarveyHoeven2021} give
\begin{equation}
 \mathsf T P_sF_s=P_tF_t\mathsf S,
 \qquad
 \|N-I\|<2.01e^{-\pi\alpha^2\theta/2}<2^{-\alpha^2\theta}.
 \label{eq:gaussian-public-facts}
\end{equation}
The first identity uses the negative Fourier sign in our definition of
$F_q$; its right-hand permutation is consequently multiplication of the
frequency by $-s$. Since $\alpha^2<p$ and $\theta>p/\alpha^4$, we have
$\alpha^2\theta>p/\alpha^2>1$. Thus $\|N-I\|<1/2$, so $N$ is invertible
by the Neumann series. Put
\[
 \mathsf S'=\mathsf S/2,\qquad
 \mathsf J'=N^{-1}/2,\qquad
 \mathsf D'=2^{-(2\alpha^2-2)}\mathsf D.
\]
The proof of Proposition~4.7(i) in the cited paper gives
\[
 \|\mathsf S'\|<\frac34,\qquad
 \|\mathsf J'\|<\frac78,\qquad
 \|\mathsf D'\|<1.
\]
Now set $A=\mathsf S'$ and $B_0=\mathsf D'\mathsf J'C$.
Because $\|C\|=1$, both are contractions. The diagonal map $\mathsf D$
is invertible, and $C\mathsf T=N\mathsf D^{-1}$, so
\[
 (\mathsf D N^{-1}C)\mathsf T
 =\mathsf D N^{-1}(N\mathsf D^{-1})=I.
\]
Applying this left inverse to the first identity in
\eqref{eq:gaussian-public-facts} gives
$\mathsf D N^{-1}CP_tF_t\mathsf S=P_sF_s$. The normalizations in
$\mathsf D'$, $\mathsf J'$, and $\mathsf S'$ have product
$2^{-(2\alpha^2-2)}2^{-1}2^{-1}=2^{-2\alpha^2}$. Hence
\[
 B_0P_tF_tA
 =2^{-2\alpha^2}\mathsf D N^{-1}CP_tF_t\mathsf S
 =2^{-2\alpha^2}P_sF_s,
\]
which is \eqref{eq:no-sort-one-dimensional}.

\paragraph{Ordered numerical maps.}
Lemmas~4.8, 4.9 and 4.12 of
\cite{HarveyHoeven2021} supply the following algorithms. Their inputs
and outputs are in increasing coordinate order, and every output is in
the disk grid. The record conversions and tape movement needed to use
these routines in our format are accounted for below. Their numerical
guarantees are
\[
\begin{array}{c|c|c}
 \text{map}&\text{error bound}&\text{time bound}\\ \hline
 \widetilde{\mathsf S}'&16p&O(tp^{3/2+\delta}\alpha)\\
 \widetilde{\mathsf J}'&3p^2/4&O(tp^{3/2+\delta}\alpha)\\
 \widetilde{\mathsf D}'&4&O(tp^{1+\delta}).
\end{array}
\]
For $C$, observe that $q_0=0$, $q_{s-1}<t$, and $q_{j+1}>q_j$,
because $t/s>1$. Thus $C$ selects an ordered subsequence. A forward scan,
with the current input coordinate stored on a work tape, copies a record
exactly when that coordinate equals the next $q_j$.

The selected indices can themselves be generated in order using only
$O(p)$ work per index. Precompute $t=hs+\rho$ with $0\leq\rho<s$.
Start with $f_0=e_0=0$, and maintain
$tj=sf_j+e_j$ with $0\leq e_j<s$. At index $j$, first use
$q_j=f_j+\mathbf1_{2e_j\geq s}$. If $j<s-1$, after copying that record
form the values for index $j+1$ by
\[
 c_j=\mathbf1_{e_j+\rho\geq s},\qquad
 e_{j+1}=e_j+\rho-c_js,\qquad
 f_{j+1}=f_j+h+c_j.
\]
Here $e_j+\rho<2s$, so $c_j$ is the only possible carry; the displayed
update gives $t(j+1)=sf_{j+1}+e_{j+1}$. All these counters have $O(p)$
bits. Elementary division computes $h,\rho$ in
$O((\log t)^2)=O(tp)$ time. The record scan and its coordinate counter
also cost $O(tp)$, so $C$ takes $O(tp)$ time with zero numerical error.

Consequently we may take
\[
 \widetilde A=\widetilde{\mathsf S}',
 \qquad
 \widetilde B_0=\widetilde{\mathsf D}'\,
                    \widetilde{\mathsf J}'\,C.
\]
Contraction of the exact maps yields errors less than $16p<p^2$ and
$4+3p^2/4<p^2$, respectively. The costs add to the claimed bound.

The original proof of Proposition~4.7(ii) implements
$B=P_s^{-1}\mathsf D'\mathsf J'CP_t$. Its two label-and-sort operations
implement $P_t$ and $P_s^{-1}$, each costing $O(tp\log t)$.
In \eqref{eq:no-sort-one-dimensional} these permutations remain in the
transform identity. The algorithms for $A$ and $B_0$ use $C$ and the
three quoted arithmetic maps, with all their original truncations.

\paragraph{Record conversion and tape movement.}
We specify how these scalar records meet the cited numerical interfaces.
Sections~2.1--2.2 of the cited paper use standard binary integers, represent
$2^{-p}(a+ib)\in\mathbb D_p$ by the exact integer pair $(a,b)$ with $p$
known, and store vectors as ordered coordinate lists. The cited
exponential and Gaussian routines return grid values in this
representation. Their computed numerator pairs are exact representations
of the approximations; the numerical errors remain those stated for the
approximated analytic maps. The source leaves the sign code, digit order,
padding and field boundaries open, so we make a concrete choice when
implementing its constructive algorithms.

At each scalar input or output boundary of a cited routine, use one sign
symbol, the ordinary binary magnitude in most-significant-bit-first order,
and a fixed separator for
each integer. Use one zero digit and the nonnegative sign for zero, put
the real field before the imaginary field, and concatenate pairs in
coordinate order. Each grid numerator magnitude has at most $p+1$ digits.
To convert one field to our $w_{\rm c}$-bit word, buffer its magnitude, right-align
it in a zero-filled block of that width, and, for a negative sign,
complement the block and add one from right to left. Copy the block from
left to right to the output. The inverse conversion reads one complete
$w_{\rm c}$-bit word, remembers its leading sign, complements and adds
one if negative, and emits the signed magnitude after removing leading
zeroes. In both directions the represented integer is unchanged.

A unary ruler of $w_{\rm c}$ cells is prepared once per call by starting
with an empty ruler and reading the binary digits of $p$ from most to least
significant, successively doubling the current unary prefix
and appending one cell when the next digit is one, then adding two cells.
This takes $O(p)$ time. Reusing the ruler, a constant
number of sweeps of the current field and work blocks performs each
conversion, including head returns and cleanup, in $O(p)$ time on fixed
tapes. Separators identify the next source field, and the ruler identifies
the next fixed-width word; the input head finishes after the consumed
field and the output head after the appended field. The inverse adapter
supplies the Gaussian input vectors, and the forward adapter formats their
outputs. Lemma~2.12 instead receives the ordinary binary integers $j,k$
and needs only the output adapter. These value-preserving conversions add
no error. On the Gaussian line lengths $s<t$, their
$O((s+t)p)=O(tp)$ work, and their $O(p)$ work per exponential output,
are included in the stated bounds.

The remaining tape movement inside the quoted algorithms also has the
claimed bound. The $\mathsf S'$ procedure uses a consecutive periodic input
window of radius $\lceil\sqrt p\rceil\alpha$ for each output. Visiting
these $O(p)$-bit records costs $O(tp\sqrt p\alpha)$ tape steps.
Remark~4.10 in the cited paper charges the extra jumps at the cyclic
boundary by the same bound.

The Neumann evaluation also keeps its computed records bounded.
Put $E=N-I$ and $k=\lceil p/(\alpha^2\theta)\rceil$. Under the present
hypotheses, Lemma~4.11 gives, for every $z\in\mathbb D_p^s$,
\[
 \widetilde E z\in\mathbb D_p^s,\qquad
 2^p\|\widetilde E z-Ez\|<p/3.
\]
Its proof establishes the rounded output bound
\[
 \|\widetilde E z\|<0.42\|z\|+(p/3)2^{-p}<1.
\]
For $u\in\mathbb D_p^s$, Lemma~4.12 rounds the real and imaginary parts
of each coordinate of $u/2$ toward zero at precision $p$, obtaining
$\widetilde v^{(0)}\in\mathbb D_p^s$. It then computes
$\widetilde v^{(j)}=\widetilde E\widetilde v^{(j-1)}$
for $1\leq j<k$. Induction therefore keeps every iterate in
$\mathbb D_p^s$. The approximation is
$\widetilde{\mathsf J}'u=\sum_{j=0}^{k-1}(-1)^j\widetilde v^{(j)}$.
Since $k\leq\alpha^2+1<p+1$, each component of any signed partial
numerator sum has magnitude at most $k2^p$ and needs
$p+\lceil\log_2(k+1)\rceil+O(1)=O(p)$ bits. Inside one $\widetilde E$
row, the $2m$ disk-grid summands of Lemma~4.11 likewise need
$p+\lceil\log_2(2m+1)\rceil+O(1)=O(p)$ accumulator bits, where
$m=\lceil\sqrt p/(2\alpha)\rceil$. These partial sums are accumulated
exactly on separate tapes; only completed iterates are fed into
$\widetilde E$.

The window of radius $m$ in one $\widetilde E$ call has tape movement
cost $O(tp(1+\sqrt p/\alpha))$, including the analogous boundary
accounting. The exact row additions cost $O(smp)$, within this bound.
The $k-1$ calls therefore spend
\[
 O\bigl((\alpha^2+1)tp(1+\sqrt p/\alpha)\bigr)
 =O(tp\sqrt p\alpha)
\]
on these scans and additions, using $2\leq\alpha<\sqrt p$.
For either Gaussian procedure, if a window
radius $w$ satisfies $s\leq2w$, then $O(1+w/s)$ traversals of the period
cost $O((s+w)p)=O(wp)$ per output. This also covers windows longer than
one period. The current iterate, next iterate, and running sum use a
fixed number of tapes holding $O(p)$-bit records. Outer accumulation
costs $O(ksp)=O(tp\alpha^2)$, within $O(tp\sqrt p\alpha)$.
All truncations are those of the cited algorithms, and the tape count
and constants remain uniform in the parameters.
\end{proof}

\subsection{Growing dimension and the padded array}

For the tensor product, apply the one-dimensional maps to one coordinate
line at a time. The chunk-swap algorithm supplies the movement of the
physical axis fields, including its cost when the number of axes grows.

\begin{lemma}[Tensor interface]
\label{lem:tensor-resampling}
Suppose every pair $(s_i,t_i)$, $1\leq i\leq d$, satisfies
Lemma~\ref{lem:no-sort-resampling} with the same $p,\alpha,\delta$.
Put $T=\prod_i t_i$, $\gamma=2d\alpha^2$, and
\[
 R=\bigotimes_iP_{s_i},\quad Q=\bigotimes_iP_{t_i},\quad
 \mathcal A=\bigotimes_i A_i,\quad
 \mathcal B_0=\bigotimes_i B_{0,i}.
\]
Then
\begin{equation}
 RF_{\boldsymbol s}
   =2^\gamma\mathcal B_0QF_{\boldsymbol t}\mathcal A,
 \qquad F_{\boldsymbol q}:=\bigotimes_iF_{q_i}.
 \label{eq:no-sort-tensor}
\end{equation}
Also write $F_{\boldsymbol q}^+=\bigotimes_iF_{q_i}^+$.
Both $\mathcal A$ and $\mathcal B_0$ are contractions. The linewise
implementations of $\mathcal A$ and $\mathcal B_0$ each have error
less than $dp^2$. The total time spent in the
one-dimensional algorithms for both maps is
$O(dTp^{3/2+\delta}\alpha)$, with a constant independent of $d$.

If each $t_i\in\{r/2,r\}$ with $r=2^\ell$, the array is kept in the
full padded box of $T$ entries, and each entry occupies $O(p)$ bits,
then exposing and restoring all the lines costs
\[
 O\bigl(Tp\,d^2(1+\ell^\tau)\bigr)
\]
using the equal-width chunk-swap bound $O(Vu^\tau)$ and the fixed-digit
movement primitive. Any fixed polynomial in $p$ of setup per line is
negligible compared with $Tp$ when $r$ grows faster than every such
polynomial and $d$ is polynomially bounded in $p$.
\end{lemma}

\begin{proof}
Tensoring \eqref{eq:no-sort-one-dimensional} proves
\eqref{eq:no-sort-tensor}. To implement $\mathcal A$, first embed its
input in the padded box by writing zero outside
$\prod_i[0,s_i)$. Expand the axes one at a time. After the axes in a set
$I$ have been expanded, the valid entries have coordinate ranges
$[0,t_i)$ for $i\in I$ and $[0,s_i)$ otherwise; every other entry is
zero. Expose the next axis $i$ as the innermost address field. On a line
whose other coordinates are valid, read its first $s_i$ values and apply
$\widetilde A_i$, then write the resulting $t_i$ values. On every other
line write $t_i$ zeros. The new valid set is the one for $I\cup\{i\}$.

For compression, start with all coordinates in $[0,t_i)$ valid.
After a set $J$ of axes has been compressed, validity means
$j_h<s_h$ for $h\in J$ and $j_h<t_h$ otherwise. For the next axis $i$,
apply $\widetilde B_{0,i}$ to each $t_i$-line whose other coordinates
are valid. Write its $s_i$ outputs followed by $t_i-s_i$ zeros.
A line invalid on another coordinate is already zero; consume it and
write a zero line without calling $\widetilde B_{0,i}$. The box therefore
keeps its full size, and the new valid set is the one for $J\cup\{i\}$.

For a fixed axis, the exact map acts independently on its valid lines
and is zero on the other lines. Taking the maximum over the line norms
shows that this axis map is a contraction; taking the maximum of the
line errors gives an error less than $p^2$ for the whole axis.
Composing the $d$ axis maps gives error less than $dp^2$.
At most $T/t_i$ lines invoke the one-dimensional algorithm for axis $i$,
so their cost is bounded by
\[
 \sum_{i=1}^d\frac{T}{t_i}
           O(t_i p^{3/2+\delta}\alpha)
   =O(dTp^{3/2+\delta}\alpha).
\]
Initial padding costs $O(Tp)$, and scanning and zeroing the remaining
lines costs $O(Tp)$ per axis. Both are within the same bound.
All calls reuse one fixed collection of work
tapes. Loop indices and the growing shape descriptor reside on tapes;
they are not extra heads.

The descriptor records which coordinate each physical address field
holds. Thus moving fields changes their positions, but not the validity
conditions above. To move an axis past its neighbor, interchange their
$\ell$-bit or $(\ell-1)$-bit address fields. Equal widths use
Lemma~\ref{lem:chunk-swap}. For unequal widths, write the longer field as
a leading bit $e$ followed by $\ell-1$ bits. If it comes first, the order
is $e\,x\,y$, with $x,y$ of equal width. Interchange $x,y$, then move $e$
between them, obtaining $y\,e\,x$. If the longer field comes second,
first move its leading bit left of the shorter field, then interchange
the two equal-width pieces. The bit move uses
Lemma~\ref{lem:elementary-streams}. Each adjacent interchange therefore
costs $O(Tp(1+\ell^\tau))$. Exposing and restoring one axis needs
$O(d)$ adjacent interchanges, hence all axes need $O(d^2)$. These are
exact permutations of the padded box, so they change neither norms nor
errors.

Since $t_i\geq r/2$, each tensor map invokes at most
$\sum_iT/t_i\leq2dT/r$ lines. For any fixed setup degree $c$, work
$O(p^c)$ per line totals $O(dTp^c/r)=o(Tp)$: after division by $Tp$,
the ratio is $O(dp^{c-1}/r)$, which tends to zero under the stated growth
assumptions. This estimate includes local descriptor copying, validity
tests per line, and preparing line boundaries.
\end{proof}

\subsection{The middle transform and the final permutation}

The factor $QF_{\boldsymbol t}$ in \eqref{eq:no-sort-tensor} can be
computed from a cyclic convolution and two pointwise phase
multiplications. The chirp identity below performs the frequency change
in $Q$ without applying $Q$ as a separate permutation of the array.
The second identity shows how the
retained source permutation $R$ cancels when computing a convolution.
Here $*$ denotes
cyclic convolution on the product of the coordinate groups and a dot
denotes coordinatewise multiplication.

\begin{lemma}[Chirp and permutation cancellation]
\label{lem:chirp-permutation}
With $R,Q,s_i,t_i$ as in Lemma~\ref{lem:tensor-resampling},
assume every $t_i$ is even. For $\boldsymbol j\in\mathbb Z^d$, define
\[
 a_{\boldsymbol j}
 =\exp\left(-\pi i\sum_{i=1}^d\frac{s_i j_i^2}{t_i}\right).
\]
Then $a$ has period $t_i$ in coordinate $i$, and
\begin{equation}
 QF_{\boldsymbol t}u
 =\overline a\mathbin{\cdot}
       \frac{a*(\overline a\mathbin{\cdot}u)}{T}.
 \label{eq:permuted-chirp}
\end{equation}
If $S=\prod_i s_i$, then
\begin{equation}
 (RF_{\boldsymbol s}^+)
   \bigl[(RF_{\boldsymbol s}u)\mathbin{\cdot}
         (RF_{\boldsymbol s}v)\bigr]
 =\frac{u*v}{S^2}.
 \label{eq:source-permutation-cancel}
\end{equation}
Complex conjugation of the input and output of an algorithm for
$RF_{\boldsymbol s}$ gives an algorithm for $RF_{\boldsymbol s}^+$.
\end{lemma}

\begin{proof}
Replacing $j_i$ by $j_i+t_i$ changes the exponent defining $a$ by
$-\pi i s_i(2j_i+t_i)$, an integral multiple of $2\pi i$.
The coefficient of $u_{\boldsymbol k}$ in the right side of
\eqref{eq:permuted-chirp} is
\[
 \frac1T\overline a_{\boldsymbol j}
               a_{\boldsymbol j-\boldsymbol k}
               \overline a_{\boldsymbol k}
 =\frac1T\exp\left(2\pi i\sum_i
                       \frac{s_i j_i k_i}{t_i}\right).
\]
This is precisely the coefficient in $QF_{\boldsymbol t}$, because
$Q$ reads frequency $(-s_i j_i)_i$ from the negative-sign transform.

For a unit $c$ modulo $q$, let $(R_cz)_j=z_{cj}$. Substituting $h=ck$
in the Fourier sum gives
\[
 (F_q^+R_cz)_j
 =\frac1q\sum_{h=0}^{q-1}z_h e^{2\pi i h c^{-1}j/q}
 =(R_c^{-1}F_q^+z)_j.
\]
Each $t_i$ is a unit modulo $s_i$, so tensoring this identity gives
$F_{\boldsymbol s}^+R=R^{-1}F_{\boldsymbol s}^+$ and hence
$RF_{\boldsymbol s}^+R=F_{\boldsymbol s}^+$. Also
$(RF_{\boldsymbol s}u)\mathbin{\cdot}(RF_{\boldsymbol s}v)
=R(F_{\boldsymbol s}u\mathbin{\cdot}F_{\boldsymbol s}v)$,
because a coordinate permutation commutes with pointwise products.
Therefore
\[
 RF_{\boldsymbol s}^+\bigl[R(F_{\boldsymbol s}u\mathbin{\cdot}
                                       F_{\boldsymbol s}v)\bigr]
 =F_{\boldsymbol s}^+(F_{\boldsymbol s}u\mathbin{\cdot}
                                       F_{\boldsymbol s}v).
\]
The normalized convolution rule is
$F_{\boldsymbol s}(u*v)=S(F_{\boldsymbol s}u)\mathbin{\cdot}
(F_{\boldsymbol s}v)$, while
$F_{\boldsymbol s}^+F_{\boldsymbol s}=I/S$. The two normalizations give
\[
 F_{\boldsymbol s}^+
   (F_{\boldsymbol s}u\mathbin{\cdot}F_{\boldsymbol s}v)
 =\frac1S F_{\boldsymbol s}^+F_{\boldsymbol s}(u*v)
 =\frac{u*v}{S^2},
\]
proving \eqref{eq:source-permutation-cancel}. Finally, conjugation
changes the negative Fourier sign to the positive one and commutes with
the real permutation $R$. It preserves both the disk grid and the norm,
so conjugating the input and output also preserves an algorithm's error
bound.
\end{proof}
